\documentclass[10pt,a4paper]{amsart}
\usepackage[margin=19mm]{geometry}
\usepackage{amsmath,amssymb,mathtools}
\usepackage[scr=rsfs,cal=euler]{mathalfa}
\usepackage{xfrac}
\usepackage[unicode,hidelinks,bookmarks=false]{hyperref}
\newcommand{\ie}{i.e.\ }
\newcommand{\cf}{cf.\ }
\newcommand{\resp}{respectively\ }
\newcommand{\Subsection}{\subsection}
\providecommand{\nobreakdash}{\nobreak\hbox{-}}
\setlength{\emergencystretch}{2em}
\multlinegap0pt
\usepackage{bm}
\usepackage{bbm}
\usepackage{dsfont}
\usepackage{xspace}
\usepackage{cancel}
\usepackage{mathtools}
\usepackage{amsthm}
\makeatletter
\@ifundefined{assumption}{\newtheorem{assumption}{Assumption}}{}
\makeatother
\newcommand{\E}{\mathbb{E}}
\title[Dynamic Decision-Making under Model Misspecification: A Stoc]{Dynamic Decision-Making under Model Misspecification: A Stochastic Stability Approach}
\author{Xinyu Dai, Daniel Chen, Yian Qian}
\date{Source: arXiv:2602.17086v2 (CC BY 4.0)}
\begin{document}
\maketitle

\noindent\emph{Source and licence.} Dynamic Decision-Making under Model Misspecification: A Stochastic Stability Approach. Version arXiv:2602.17086v2 (CC BY 4.0), distributed under the Creative Commons Attribution 4.0 International licence, as recorded in the arXiv licence metadata for this exact version and shown on the abstract page https://arxiv.org/abs/2602.17086v2. Full rights notice: licenses/arxiv-2602.17086-CC-BY-4.0.txt.\par\smallskip

\noindent\textit{Editorial scope.} Counted unit: the Assumption whose source label is \texttt{as:two-arm} in the article (source0.tex lines 1641--1650, source label \texttt{as:two-arm}), delivered together with the article's own complete original proof of it (source0.tex lines 1652--1717, one \texttt{proof} environment of 66 lines). No mathematics is written, repaired or re-derived: statement and proof are the article's own text line for line. Required context retained uncounted: eq:filtration (lines 170--173). Every definition, display and label of those lines is retained unchanged. Omitted from this package and not redistributed: 1--169, 174--1640, 1651--1651, 1718--2554. Nothing inside a retained excerpt is omitted or altered: every retained body is delivered exactly as the article's lines are written, so no omission receipt of any kind accompanies this package. Citations: the retained excerpts carry no citation command at all, so no bibliography entry of the article is transcribed and no reference is redistributed. Reference adaptation: none was needed and none was made; every reference inside the delivered unit resolves inside it, so no unresolved reference or citation is produced. Printed slips kept literal: no printed word, symbol, hypothesis or number of the counted statement or of its proof was altered. Layout and class: the article's own class is \texttt{article} with packages including amsmath, amsthm, amssymb, newpxtext, newpxmath, graphicx, subcaption, setspace, hyperref, float, listings, color, tikz, pgfplots; the wrapper is the lane's pinned \texttt{amsart} editorial wrapper at 10pt on A4, which restates the theorem-like environments the retained text needs and adds the article's own macro definitions that the retained text needs (\texttt{\textbackslash E}), with extra wrapper packages (bm, bbm, dsfont, xspace, cancel, mathtools). The delivered numbering is the unit's own. Assets: no retained span contains a figure environment, a graphics-inclusion command, a PDF-page inclusion, a table or a TikZ picture; no image file is copied into this package, the package declares no asset dependency and no image file is redistributed with it. Licence: version arXiv:2602.17086v2 of this article is distributed under the Creative Commons Attribution 4.0 International licence, as recorded in the arXiv licence metadata for this exact version and shown on the abstract page https://arxiv.org/abs/2602.17086v2. A full rights notice is delivered at licenses/arxiv-2602.17086-CC-BY-4.0.txt. Characters: every retained excerpt is pure ASCII, so the delivered engine, XeLaTeX with its default Latin Modern text font, reports no missing character.
\begin{equation}
\mathcal{F}_t = \sigma(A_0, R_0, A_1, R_1, \ldots, R_{t-1}, A_t).
\label{eq:filtration}
\end{equation}

\begin{assumption}
We assume the following: 
\begin{enumerate}
    \item initial log-likelihood does not prefer any arm, i.e., $S_0 = 0$;
    \item Under the true reward distribution given action $i$ $g_i$, the log-likelihood
    \begin{align*}
        \E_{g_i} \left[ \left( \log \frac{f_{\nu_i}(r)} {f_{\gamma_i}(r)} \right)^2 \right] < \infty;
    \end{align*}
\end{enumerate}\label{as:two-arm}
\end{assumption}

\begin{proof}
Assume both models agree on the optimal action:
\[
\phi(\nu) = \phi(\gamma) = 1.
\]
Thus, Thompson Sampling always selects action 1: \( A_t = 1 \) for all \( t \). 
Let \( R_t \) denote the reward observed at time \( t \) upon performing action 1, and define the log-likelihood ratio process as
\[
S_t := \sum_{s=1}^t \log \frac{f_{\nu_1}(R_s)}{f_{\gamma_1}(R_s)}.
\]
Since rewards are assumed to be independent and only arm 1 is being played, the collection of random variables $(Z_s)_{s \in \mathbb N}$ defined as
\[
Z_s := \log \frac{f_{\nu_1}(R_s)}{f_{\gamma_1}(R_s)}
\]
are i.i.d.~and integrable by assumption.
Denote their expectations by
\[
\Delta_1 := \mathbb{E}_{g_1} \left[ \log \frac{f_{\nu_1}(R)}{f_{\gamma_1}(R)} \right].
\]

Recall the filtration $\mathbb F = (\mathcal F_t)_{t \in \mathbb N}$ generated by previous actions and rewards in \eqref{eq:filtration}.
Then \( \{S_t\} \) is $\mathbb F$-adapted and satisfies
\[
\mathbb{E}[S_{t+1} \mid \mathcal{F}_t] = S_t + \Delta_1.
\]
We analyze three cases based on the sign of \( \Delta_1 \):

\paragraph{\underline{Case (i): $\Delta_1 > 0$}}

The process \( S_t \) is a \emph{submartingale} under \( \mathbb{P} \) since
\[
\mathbb{E}[S_{t+1} \mid \mathcal{F}_t] = S_t + \Delta_1 \geq S_t.
\]
Since \( S_t \) is a sum of i.i.d.~variables with positive mean, the strong law of large numbers implies
\[
\lim_{t \to \infty} \frac 1 t S_t = \Delta \quad \text{$\mathbb P$-a.s.},
\]
and consequently, $S_t$ diverges a.s.
The posterior belief evolves as:
\[
\pi_t = \frac{1}{1 + \frac{1 - \pi_0}{\pi_0} e^{-S_t}} \to 1 \quad \text{a.s. as $t \to \infty$.}
\]

\paragraph{\underline{Case (ii): $\Delta_1 = 0$}}

In this case, the log-odds process $S$ is a symmetric random walk.
So, for any $z \in (0,1)$, we have
\begin{align*}
    \Pr(\pi_t \leq z) = \Pr \left( S_t \leq - \log \left( \frac 1 z - 1 \right) \right) = \Pr \left( \frac{S_t}{\sqrt t} \leq \frac {- \log \left( \frac 1 z - 1 \right)} {\sqrt t} \right).
\end{align*}
Now, by the independence of the reward and the assumed second-moment bound, we know that a central limit theorem holds.
Furthermore, by Slutsky's theorem, we have
\begin{align*}
    \frac{S_t}{\sqrt t} + \frac {- \log \left( \frac 1 z - 1 \right)} {\sqrt t} \stackrel{\text{(d)}}{\longrightarrow} \mathcal N(0, \sigma^2).
\end{align*}
Since $0$ is a point of continuity for the normal distribution, we take the limit on both sides and conclude that
\begin{align*}
    \lim_{t \to \infty} \Pr(\pi_t \leq z) = \Pr(\mathcal N(0, \sigma^2) \leq 0) = \frac 1 2.
\end{align*}
By the fact that $z$ is chosen arbitrarily and $0 \leq \pi_t \leq 1$, this shows that $\pi_t$ converges weakly to a Bernoulli with parameter $1/2$.

\paragraph{\underline{Case (iii): $\Delta_1 < 0$}}

This case follows from Case (i) by symmetry, and the proof is complete. 

\end{proof}

\end{document}
